\documentclass[a4paper,addpoints,14pt]{exam}
%\documentclass[answers,addpoints,14pt]{exam}
\unframedsolutions

\usepackage{array,amsmath,fancybox,amsfonts,graphicx,amsthm,amssymb,wasysym}
\usepackage{latexsym,slashed,anysize,subfigure,calc}
\usepackage[spanish]{babel}
\usepackage[utf8]{inputenc}
%\usepackage{fancyhdr,float,hyperref}
%\usepackage{lettrine}

%\usepackage{enumitem}

%entorno para química: texto y math mode
%\newcommand*\chem[1]{\textup{#1}}
\newcommand*\chem[1]{\ensuremath{\mathrm{#1}}}

%tamaños fuentes extra: 8,9,14,17,20
	\usepackage{extsizes}

%display style in long fractions
\newcommand\ddfrac[2]{\frac{\displaystyle #1}{\displaystyle #2}}

%tabla generada
\usepackage[table,xcdraw]{xcolor}

%lewis y química
\usepackage{chemformula,elements}
%chemformula, provides \ch{} with the !(<below>)(<formula>) syntax and \chlewis;
%elements, provides \elconf and \writeelconf.
%ex:
%\ch{
%  !(\elconf{Li})( "\chlewis{180.}{Li}" ) +
%  !(\elconf{F})( "\chlewis{0.90:180:270:}{F}" )
%   ->
%  !(\writeelconf{2})( Li+ ) +
%  !(\writeelconf{2,2+6})( "\chlewis{0:90:180:270:}{F}" {}- )
%}

%multicol
\usepackage{multicol}
  

%símbolo bonito para qed
\usepackage{amsthm}
	%\newcommand{\mc}[3]{\multicolumn{#1}{#2}{#3}}
\newcommand{\halmos}{\vspace*{-6mm}
	
					\begin{flushright}
						\rule{1mm}{2.5mm} 
					\end{flushright}
					\vspace*{-4mm}
						
					  }
\renewcommand{\qed}{\halmos}

 
%página blanco
\usepackage{afterpage}
%uso \afterpage{\blankpage}

\newcommand\blankpage{%
    \null
    \thispagestyle{empty}%
    \addtocounter{page}{0}%
    \newpage}





 %caja bonita
  \usepackage[many]{tcolorbox}
\definecolor{titlebg}{RGB}{80,100,72}
\newtcolorbox{learn}{
  breakable,
  enhanced,
  arc=0pt,
  outer arc=0pt,
  colframe=titlebg,
  colback=titlebg!7,
  overlay unbroken and first={
    \node[
      draw=titlebg,
      fill=titlebg,
      rotate=270,
      anchor=north west,
      text=white,
      font=\bfseries
    ]
    at (frame.north west)  
    {DATOS};
  }
}


\newtcolorbox{formulacion}{
  breakable,
  enhanced,
  arc=0pt,
  outer arc=0pt,
  colframe=titlebg,
  colback=titlebg!7,
  overlay unbroken and first={
    \node[
      draw=titlebg,
      fill=titlebg,
      rotate=270,
      anchor=north west,
      text=white,
      font=
    ]
    at (frame.north west)  
    {\hspace*{3ex}};
  }
}

\newtcolorbox{problemas}{
  breakable,
  enhanced,
  arc=0pt,
  outer arc=0pt,
  colframe=titlebg,
  colback=titlebg!7,
  overlay unbroken and first={
    \node[
      draw=titlebg,
      fill=titlebg,
      rotate=270,
      anchor=north west,
      text=white,
      font=
    ]
    at (frame.north west)  
    {AVISO};
  }
}
	
\usepackage{anysize}
	%\marginsize{left}{right}{top}{bottom}
    \marginsize{1cm}{.8cm}{1cm}{1cm}
    
\usepackage[spanish]{babel}
\usepackage[utf8]{inputenc}

\usepackage{sidecap}
\usepackage{mwe,float,afterpage}

%redefiniciones
\pointpoints{punto}{puntos}
\bonuspointpoints{punto}{puntos}
\hqword{Pregunta} %substitutes text for \Question"
%\vpgword{Pág.}%substitutestext for \Page"
\hpword{Máx.}%substitutes text for \Points"
\hsword{Nota}%substitutes text for \Score"
\htword{Total}%substitutes text for \Total:"

\chqword{Pregunta}
\chpword{Máx.}
\chsword{Nota}
\chtword{Total}
\chbpword{Puntos \textit{bonus}}

\renewcommand{\solutiontitle}{\noindent\textbf{Solución:}\par\noindent}


%fuentes
	%\usepackage[T1]{fontenc}
	%\usepackage{fourier}

	\usepackage[T1]{fontenc}
	\usepackage{librebaskerville}


\usepackage{amssymb,amsmath}
\usepackage{graphicx,graphics}
\usepackage{color}
	%	\colorgrids %grids con color
	%	\definecolor{GridColor}{gray}{0.8}
		%grid
%			\setlength{\gridsize}{5mm}
%   			 \setlength{\gridlinewidth}{0.1pt}

%fracciones en texto más elaboradas
	\usepackage{xfrac}

%mejora en la exposición de guarismos
	\usepackage{numprint}
 
%caption
\usepackage[font=small,labelfont=bf,
labelsep=endash]{caption}
	
%soluciones sombreadas	
\shadedsolutions


\checkboxchar{$\Box$}
\checkedchar{$\blacksquare$}

\usepackage{everypage,draftwatermark} %Marca de agua
\SetWatermarkText{\hspace*{-1.5cm} \parbox{\textwidth}{\centering  123}}
\SetWatermarkScale{3}
\SetWatermarkLightness{0.9}


%cabecera
\header{\includegraphics[scale=.22]{a.png}}       
      {}
       {$f_{ord}$\textbf{B1-{\small Dist.}}\;{\small I-II-III}\\14\textbf{mayo}2024 }       
%pie de página

\cfoot{}
   %enumera en todas menos en la primera
   %\runningfootrule

   %enumera en todas
   %\footrule
   %\footer{}{Página \thepage\ de \numpages}{}
%   \footer{}
%	{Página \thepage\ de \numpages}
%	{\iflastpage{{\footnotesize \textit{Fin del examen}}}{{\footnotesize \textit{Da la vuelta a la hoja}\ldots}}}


		
\begin{document}


\fbox{\parbox{0.945\linewidth}{\textbf{Nombre y apellidos}:\vspace*{20mm}}}


\begin{learn}
\begin{multicols}{2}
{$\triangleright\; M(\chem{O})\simeq16\;g\,mol^{-1}$\newline
$\triangleright\; M(\chem{C})\simeq12\;g\,mol^{-1},$\newline
$\triangleright\; M(\chem{N})\simeq14\;g\,mol^{-1},$\newline
$\triangleright\; M(\chem{H})\simeq1\;g\,mol^{-1},$\newline
$\triangleright\; M(\chem{Na})\simeq 23\;g\,mol^{-1},$\newline
$\triangleright\; M(\chem{Cl})\simeq 36\;g\,mol^{-1},$\newline
$\triangleright\; M(\chem{S})\simeq32\;g\,mol^{-1},$\newline
$\triangleright\; M(\chem{Mn})\simeq55\;g\,mol^{-1},$\newline
$\triangleright\; M(\chem{Mg})\simeq 24\;g\,mol^{-1},$\newline
$\triangleright\; R\simeq 0.082\;atm\,L\,mol^{-1}\,K^{-1}$
}
\end{multicols}
\end{learn}	

\vspace{2ex}

\begin{problemas}{Esta prueba tiene \numquestions\ preguntas y un total de 
\numpoints\  puntos. Recuerda que un examen es una \textit{comunicación}: debes invocar las leyes en que te basas para resolver los ejercicios así como \textit{indicar brevemente los pasos y razonamientos} que sigues para obtener los resultados.}
\end{problemas}
\vspace{9ex}

\qformat{\textbf{Pregunta \thequestion}\quad \hfill} 
\boxedpoints

\begin{questions}

\question En un matraz de $5.0\;L$ se introducen $300\;g$ de oxígeno molecular y $150\;g$ de nitrógeno molecular. Si la temperatura es de $30^\circ C$ y asumimos que ambos gases se comportan idealmente, responde \textit{razonadamente} a las siguientes cuestiones:
	\begin{parts}
		\part[1] Calcula la presión total de la mezcla.
		\part[1] Comprueba que se cumple la \textit{ley de Dalton}.
	\end{parts}
\qed


\question Considera los elementos desconocidos $^{8}\alpha$, $^{18}\beta$, $^{34}\delta$ y $^{38}\epsilon.$ 
    	\begin{parts}
    	\setlength\itemsep{0.3ex}
    		\part[1] Escribe sus {configuraciones electrónicas}. 
    		\part[1] {Describe las propiedades} (estructura, solubilidad, conducción eléctrica y estado físico más probable a temperatura ambiente) que presentará el compuesto del apartado b).
    		\part[1] Justifica la existencia y el tipo de enlace de la molécula $\alpha_2$ .
    	\end{parts}
		\qed



\question   A $100^\circ C$ y $3\;atm$  se hacen reaccionar $9\;g$ de dióxido de manganeso sólido con $60\;mL$ de una disolución $3\;M$ de ácido clorhídrico , obteniéndose $4.8\;g$ de dicloruro de manganeso sólido además de cloro molecular gaseoso y agua líquida.
    	\begin{parts}
    		\part[1] Ajusta la reacción e identifica el reactivo limitante.
    		\part[1] Determina el rendimiento de la reacción.
    	\end{parts}
  	\qed

%\question[2] Las entalpías de formación del metanol líquido, el dióxido de carbono gaseoso y el agua líquida a presión de una atmósfera son, respectivamente, $-239\;kJ\,mol^{-1}$; $-393.5\;kJ\,mol^{-1}$ y $-285,6\;kJ\,mol^{-1}$.
%	\begin{parts}
%		\part[1] Ajusta la ecuación de combustión de $1\;mol$ de metanol y calcula la entalpía estándar molar de reacción.
%		\part[1] Determina el calor involucrado en la combustión de $1\;kg$ de metanol. ¿Es un calor \textit{absorbido} del entorno o \textit{cedido} al entorno?
%	\end{parts}
%\vspace{1ex}
%	
%\qed
%%
%%%\afterpage{\blankpage}
%%\newpage
%%\vspace*{-7ex}
%%
%%\begin{center}
%%       \fbox{\footnotesize HOJA DE RESPUESTAS A LOS PROBLEMAS}
%%\end{center}

	\question El adjetivo supersónico aplicado a balas refiere a un tipo de munición que es capaz de alcanzar la velocidad del sonido en el aire, $c_s\simeq 340\;ms^{-1}$. 
		\begin{parts}
			\part[1] Describe un {\small SR} apropiado para estudiar el problema y obtén las {ecuaciones cinemáticas vectoriales} de una bala supersónica disparada desde el cañón de un rifle que se encuentra a $1.7\;m$ del suelo y cuya inclinación respecto de la horizontal son $30^\circ$.
			\part[1] Determina el {tiempo de vuelo} del proyectil.
			\part[1] Determina la {altura máxima} sobre el suelo que se alcanza durante su trayectoria.
		\end{parts}
\qed



\question 
\begin{minipage}{10cm}
	Un objeto que puede considerarse puntual está situado sobre una rampa inclinada $25^\circ$ con la horizontal tal y como se muestra en la figura. El rozamiento de la caja con la rampa puede modelizarse a partir de un  coeficiente $\mu=0.2$.

\end{minipage}\hfill
\begin{minipage}{7.1cm}
\vspace{-2ex}

	\begin{figure}[H]
		\centering
		\def\svgwidth{1.05\linewidth}
		\input{plano1.pdf_tex}
	\end{figure}
\end{minipage}
\vspace{3ex}

	
	\begin{parts}
		\part[1] {Representa} con un diagrama de cuerpo libre todas las fuerzas que actúan sobre el objeto.
	
		\part[1] Calcula la {aceleración} con la que se mueve el objeto.
		\part[1] Calcula la {velocidad} con la que llega al suelo.
	
\vspace{1ex}

		\part[2] Si la rampa fuera móvil y pudiera ajustarse la pendiente a voluntad, calcula el {valor máximo} que puede tomar el ángulo de la pendiente si se quiere que el objeto \textit{permanezca en reposo}.
	\end{parts}
\qed
	
	

\end{questions}	


\end{document}






